\section*{Ejercicio 3}
\noindent Halle y sombree el área, usando integrales en curvas polares, de la región encerrada por la curva \(r = -8\cos(\theta)\) (circunferencia).

\begin{gather*}
  A / 2 = \frac{1}{2} \int_{0}^{\pi} \left(-8\cos(\theta)\right)^2 d\theta %
    = 32 \int_{0}^{\pi} \cos^2(\theta) \, d\theta %
    = 16 \int_{0}^{\pi} 1 + \cos(2\theta) \, d\theta %
    = 16 \left(\theta + \frac{\sin(2\theta)}{2}\right)\bigg|_{0}^{\pi} %
    = 16 \pi \\
  A = 32 \pi
\end{gather*}

\noindent El área de la región encerrada por la curva \(r = -8\cos(\theta)\) es \(32 \pi\) unidades cuadradas.

\vspace{5cm}